%============================================================
% APPENDIX — CODE LIBRARY
%   One tested C++14 header per component in code/lib/, random generators in code/gen/,
%   one brute-force test per component in code/tests/, run by tools/test_library.py.
%   Every output shown in this appendix was produced by running the library.
%============================================================
\chapter{Code Library}
\label{ch:code-library}

% ---- local helpers (prefixed cl to avoid clashes with chapter macros) ----
\providecommand{\clfile}[1]{\texttt{\detokenize{#1}}}
% a conditional page break: start a new page if fewer than ~8 lines are left
\providecommand{\clneed}{\par\vspace{0pt plus 8\baselineskip}\penalty-200\vspace{0pt plus -8\baselineskip}}
\providecommand{\cltest}[3]{%
  \par\smallskip\noindent\textbf{\textcolor{successgreen}{Tested by}} \clfile{#1}
  \hfill\texttt{\textcolor{successgreen}{OK #3}}\par\nobreak\noindent #2\par}

% ---- the running example L: 6 vertices, 0-1 0-2 1-3 2-3 3-4, vertex 5 isolated ----
\tikzset{
    clv0/.style={gv}, clv1/.style={gv}, clv2/.style={gv}, clv3/.style={gv}, clv4/.style={gv}, clv5/.style={gv},
    cle01/.style={ge}, cle02/.style={ge}, cle13/.style={ge}, cle23/.style={ge}, cle34/.style={ge},
    pics/clL/.style={code={
        \node[clv0] (-0) at (0,0)      {0};
        \node[clv1] (-1) at (1.3,0.8)  {1};
        \node[clv2] (-2) at (1.3,-0.8) {2};
        \node[clv3] (-3) at (2.6,0)    {3};
        \node[clv4] (-4) at (3.9,0)    {4};
        \node[clv5] (-5) at (5.1,0)    {5};
        \draw[cle01] (-0) -- (-1); \draw[cle02] (-0) -- (-2); \draw[cle13] (-1) -- (-3);
        \draw[cle23] (-2) -- (-3); \draw[cle34] (-3) -- (-4);
    }},
    % weighted sample W (generated by: python code/gen/gen.py connected 5 6 7 --weights 1 9 --base 0)
    clw0/.style={gv}, clw1/.style={gv}, clw2/.style={gv}, clw3/.style={gv}, clw4/.style={gv},
    clf01/.style={ge}, clf04/.style={ge}, clf41/.style={ge}, clf43/.style={ge}, clf30/.style={ge}, clf42/.style={ge},
    pics/clW/.style={code={
        \node[clw0] (-0) at (0,0)      {0};
        \node[clw1] (-1) at (1.5,1.2)  {1};
        \node[clw4] (-4) at (2.6,0)    {4};
        \node[clw3] (-3) at (1.3,-1.2) {3};
        \node[clw2] (-2) at (4.2,0)    {2};
        \draw[clf01] (-0) -- node[gw] {5} (-1);
        \draw[clf04] (-0) -- node[gw, pos=0.3] {9} (-4);
        \draw[clf41] (-4) -- node[gw] {3} (-1);
        \draw[clf43] (-4) -- node[gw] {2} (-3);
        \draw[clf30] (-3) -- node[gw] {4} (-0);
        \draw[clf42] (-4) -- node[gw] {6} (-2);
    }}
}

\begin{tcolorbox}[enhanced, colback=white, colframe=secondaryblue, boxrule=1pt, arc=2mm, left=10pt, right=10pt]
\textbf{\large\textcolor{primaryblue}{Twenty small headers, each one tested against a brute force on thousands of random graphs. Whenever a chapter needs a BFS, a DSU or a Dijkstra, it is this code --- one source of truth, so a bug fixed here is fixed everywhere.}}
\end{tcolorbox}

\chaptercontents

%============================================================
\clneed
\section{How to Use the Library}
\label{sec:cl-howto}

The library lives in the book's \clfile{code/} folder. Every component is \textbf{one self-contained header}: it includes \clfile{<bits/stdc++.h>}, opens \texttt{namespace std} (as contest code does) and defines one small \texttt{struct} or a couple of functions. Nothing depends on anything else, so in a contest you paste exactly the header you need above your \texttt{main()} --- or several of them: they never clash, and the only shared type, \texttt{Edge}, is guarded by \texttt{\#ifndef GRAPHLIB\_EDGE}.

\begin{diagrambox}
\begin{tikzpicture}[
    fb/.style={rectangle, rounded corners=3pt, draw=#1, fill=#1!8, line width=1pt, align=left,
               font=\scriptsize, inner sep=4pt, text width=3.5cm},
    hd/.style={font=\small\bfseries\ttfamily, anchor=south west}]
    \node[fb=primaryblue] (lib) at (0,0) {\textbf{code/lib/} (20 headers)\\graph, dfs\_iterative, dfs\_recursive, bfs, multi\_bfs, grid\_bfs, components, bipartite, cycle, toposort, functional\_graph, dsu, weighted\_dsu, zero\_one\_bfs, dijkstra, bellman\_ford, floyd\_warshall, kruskal, prim, lca};
    \node[fb=accentorange] (gen) at (0,-3.9) {\textbf{code/gen/}\\gen.hpp (C++) and gen.py (Python): seeded random trees, graphs, DAGs, functional graphs, grids, weights};
    \node[fb=successgreen] (tests) at (5.0,-1.9) {\textbf{code/tests/}\\test\_<component>.cpp: the header vs.\ a brute force on thousands of random graphs; check.hpp holds the brute forces};
    \node[fb=recursionred] (tl) at (10.0,-0.9) {\textbf{tools/test\_library.py}\\compiles and runs every test, prints a summary, exit code 1 on failure};
    \node[fb=softpurple] (st) at (10.0,-3.2) {\textbf{tools/stress.py}\\solution vs.\ brute vs.\ generator for YOUR problem (Appendix~\ref{ch:testing})};
    \draw[fcarrow] (lib.east) -- node[lbl, above] {includes} (tests.north west);
    \draw[fcarrow] (gen.east) -- node[lbl, sloped, below] {inputs} (tests.south west);
    \draw[fcarrow] (tests.east) -- (tl.west);
    \draw[fcarrow, dashed] (gen.south east) to[bend right=12] node[lbl, below] {gen.py} (st.south west);
\end{tikzpicture}
\figcap{The four parts of the library. Headers are what you copy; generators and tests are what make the headers trustworthy.}
\end{diagrambox}

\subsection*{Conventions}

\begin{itemize}[leftmargin=*]
    \item \textbf{Vertices are 0-indexed}: $0,\dots,n-1$, everywhere. Inputs in most judges are 1-indexed; convert \emph{once}, while reading (subtract 1), and never think about it again. \texttt{read\_graph(cin, n, m, directed, 1)} does exactly that.
    \item \textbf{Adjacency lists are plain} \texttt{vector<vector<int>>}; weighted ones are \texttt{vector<vector<pair<int,long long>>>} with pairs (neighbour, weight). An undirected edge is stored in both lists. Every header takes these plain types, so you can build them however you like.
    \item \textbf{Edge lists} use \texttt{pair<int,int>} or \texttt{Edge\{u, v, w\}} (with \texttt{long long w}).
    \item \textbf{Distances} are \texttt{int} with $-1$ for ``unreachable'' in the unweighted searches, and \texttt{long long} with \texttt{INF = LLONG\_MAX/4} in the weighted ones (so that \texttt{INF + w} never overflows).
    \item \textbf{Nothing recurses deeply.} Every algorithm that could recurse $n$ levels (DFS, cycle finding, topological sort, components, functional graphs) is iterative. The two exceptions are documented: \clfile{dfs_recursive.hpp} (the textbook form, with a depth warning) and \texttt{WeightedDSU::find} (depth at most $\log_2 n$).
    \item \textbf{Multi-test inputs.} Every \texttt{run()} / \texttt{build()} starts with \texttt{assign(...)} on all its arrays, so calling it again for the next test case is safe. DSUs are created per test: \texttt{DSU d(n);} inside the test loop.
\end{itemize}

\begin{diagrambox}
\begin{tikzpicture}[x=1cm, y=1cm]
    \node[draw=darkgray, fill=lightgray!60, rounded corners=2pt, font=\ttfamily\small, align=left, inner sep=4pt] (in) at (0,0) {4 3\\1 2\\2 3\\4 1};
    \node[font=\scriptsize\color{darkgray}] at (0,1.25) {judge input (1-indexed)};
    \draw[fcarrow] (1.0,0) -- node[lbl, above] {$u-1,\ v-1$} (3.0,0);
    \begin{scope}[xshift=4.3cm]
        \node[gv] (a) at (0,0.7) {0}; \node[gv] (b) at (1.3,0.7) {1}; \node[gv] (c) at (1.3,-0.7) {2}; \node[gv] (d) at (0,-0.7) {3};
        \draw[ge] (a) -- (b); \draw[ge] (b) -- (c); \draw[ge] (d) -- (a);
        \node[font=\scriptsize\color{darkgray}] at (0.65,1.45) {inside the program};
    \end{scope}
    \draw[fcarrow] (6.3,0) -- node[lbl, above] {$+1$} (7.8,0);
    \node[font=\scriptsize, align=left, anchor=west] at (7.9,0) {print answers\\1-indexed again};
\end{tikzpicture}
\figcap{Convert at the border of the program, in both directions. The library never sees a 1-indexed vertex.}
\end{diagrambox}

A typical contest file therefore looks like this (the header text is pasted where the comment is):

\begin{lstlisting}[style=snippet]
// paste code/lib/graph.hpp and code/lib/bfs.hpp here
int main() {
    int t; cin >> t;
    while (t--) {                                  // multi-test: everything is rebuilt per test
        int n, m; cin >> n >> m;
        Graph g = read_graph(cin, n, m, false, 1); // 1-indexed input -> vertices 0..n-1
        BFS b; b.run(g.adj, 0);
        for (int v = 0; v < n; v++) cout << b.dist[v] << " \n"[v == n - 1];
    }
}
\end{lstlisting}

\subsection*{The components at a glance}

\begin{table}[H]
\centering\small
\begin{tabularx}{\linewidth}{@{}>{\ttfamily\raggedright\arraybackslash}p{3.3cm} >{\raggedright\arraybackslash}X >{\raggedright\arraybackslash}p{2.4cm} r@{}}
\toprule
\textrm{\textbf{header}} & \textbf{what you get} & \textbf{time} & \textbf{cases} \\
\midrule
graph.hpp & \texttt{Graph}, \texttt{WGraph}, \texttt{read\_graph}, \texttt{CSR} & $\bigO{n+m}$ & 3000 \\
dfs\_iterative.hpp & \texttt{IterDFS}: order, parent, tin/tout, no recursion & $\bigO{n+m}$ & 4001 \\
dfs\_recursive.hpp & \texttt{RecDFS}: the same, recursive (depth note) & $\bigO{n+m}$ & 4000 \\
bfs.hpp & \texttt{BFS}: dist, par, path & $\bigO{n+m}$ & 5000 \\
multi\_bfs.hpp & \texttt{MultiBFS}: nearest source & $\bigO{n+m}$ & 5000 \\
grid\_bfs.hpp & \texttt{GridBFS}: 4/8 moves, walls, path & $\bigO{RC}$ & 3000 \\
components.hpp & \texttt{Components}: ids, sizes & $\bigO{n+m}$ & 5000 \\
bipartite.hpp & \texttt{Bipartite}: colouring or odd cycle & $\bigO{n+m}$ & 5000 \\
cycle.hpp & a cycle, undirected (multigraph) / directed & $\bigO{n+m}$ & 4000 \\
toposort.hpp & \texttt{topo\_kahn}, \texttt{topo\_dfs} (empty on a cycle) & $\bigO{n+m}$ & 5000 \\
functional\_graph.hpp & cycles, cycle id, distance to cycle, entry & $\bigO{n}$ & 5000 \\
dsu.hpp & \texttt{DSU}: union by size + path compression & $\bigO{\alpha(n)}$ & 3000 \\
weighted\_dsu.hpp & \texttt{WeightedDSU}: potentials, contradictions & $\bigO{\log n}$ & 3000 \\
zero\_one\_bfs.hpp & \texttt{ZeroOneBFS}: weights 0/1 with a deque & $\bigO{n+m}$ & 4000 \\
dijkstra.hpp & \texttt{Dijkstra}: lazy heap, par, path & $\bigO{m\log m}$ & 5000 \\
bellman\_ford.hpp & \texttt{BellmanFord} ($-\infty$ marks), \texttt{find\_negative\_cycle} & $\bigO{nm}$ & 5000 \\
floyd\_warshall.hpp & all pairs, negative-cycle test & $\bigO{n^3}$ & 5000 \\
kruskal.hpp & minimum spanning forest, chosen edges & $\bigO{m\log m}$ & 4000 \\
prim.hpp & minimum spanning forest, parents & $\bigO{m\log m}$ & 4000 \\
lca.hpp & \texttt{LCA}: binary lifting, $k$-th ancestor, dist & $\bigO{\log n}$/query & 3000 \\
\bottomrule
\end{tabularx}
\end{table}

The last column is the number of random cases each test checks against its brute force; the full run (\texttt{python tools/test\_library.py}) ends with
\begin{lstlisting}[style=trace]
20/20 library tests passed.
\end{lstlisting}

\subsection*{The running example}

Most usage examples below run on the same small graph $L$, so their outputs can be compared:

\begin{diagrambox}
\begin{tikzpicture}
    \pic (L) {clL};
    \node[draw=darkgray, fill=lightgray!60, rounded corners=2pt, font=\ttfamily\small, align=left, inner sep=4pt, anchor=west] at (6.3,0) {Graph g(6);\\g.add\_edge(0, 1); g.add\_edge(0, 2);\\g.add\_edge(1, 3); g.add\_edge(2, 3);\\g.add\_edge(3, 4);  // 5 isolated};
\end{tikzpicture}
\figcap{$L$: $n=6$, five undirected edges, a 4-cycle $0{-}1{-}3{-}2$ with a tail to $4$, and the isolated vertex $5$ --- a disconnected graph on purpose.}
\end{diagrambox}

%============================================================
\clneed
\section{Graph Storage}
\label{sec:cl-graph}

\textbf{What it does.} \texttt{Graph} is the adjacency list every other header consumes: \texttt{add\_arc(u,v)} for a directed arc, \texttt{add\_edge(u,v)} for an undirected edge (stored twice). \texttt{WGraph} is the same with weights. \texttt{read\_graph} / \texttt{read\_wgraph} read $m$ edges and shift a 1-indexed input to 0-indexed. \texttt{CSR} (\emph{compressed sparse row}) stores the same lists in two flat arrays: the neighbours of $u$ are \texttt{to[start[u]]} \ldots{} \texttt{to[start[u+1]-1]}. It is built by counting degrees, taking prefix sums and placing every edge in input order; it uses less memory than $n$ separate vectors and is friendlier to the cache when $n$ reaches $10^6$.

\textbf{Complexity.} $\bigO{n+m}$ time and memory for all of them.

\begin{diagrambox}
\begin{tikzpicture}[x=1cm, y=1cm]
    \begin{scope}[scale=0.85, transform shape]
        \pic (L) {clL};
    \end{scope}
    \node[note, anchor=west, align=left] at (5.2,0) {adj[0] = 1 2\\adj[1] = 0 3\\adj[2] = 0 3\\adj[3] = 1 2 4\\adj[4] = 3\\adj[5] = (empty)};
    \begin{scope}[yshift=-2.2cm, xshift=0.6cm, x=0.8cm]
        \node[rowlabel] at (-0.6,0) {start};
        \foreach \v [count=\i from 0] in {0,2,4,6,9,10,10} \node[cell, minimum width=7mm] (s\i) at (\i,0) {\v};
        \foreach \i in {0,...,6} \node[font=\tiny\color{darkgray}] at (\i,0.5) {\i};
        \node[rowlabel] at (-0.6,-1.4) {to};
        \foreach \v/\st [count=\i from 0] in {1/basecell,2/basecell,0/cell,3/cell,0/cell,3/cell,1/compcell,2/compcell,4/compcell,3/cell}
            \node[\st, minimum width=7mm] (t\i) at (\i,-1.4) {\v};
        \foreach \i in {0,...,9} \node[font=\tiny\color{darkgray}] at (\i,-1.9) {\i};
        \draw[dep] (s3.south) -- (t6.north);
        \draw[dep] (s4.south) -- (t9.north);
        \node[note, anchor=west, align=left] at (10.4,-0.7) {neighbours of 3:\\to[6..8] = 1 2 4\\(start[3] = 6,\\start[4] = 9)};
    \end{scope}
\end{tikzpicture}
\figcap{$L$ as adjacency lists and as CSR (real output of the example below). The \texttt{to} array has $2m=10$ entries; the two shaded blocks are the neighbours of 0 and of 3.}
\end{diagrambox}

\begin{lstlisting}[style=snippet]
vector<pair<int,int>> E = {{0,1}, {0,2}, {1,3}, {2,3}, {3,4}};
Graph g(6);
for (auto& e : E) g.add_edge(e.first, e.second);
CSR c(6, E, false);                         // false = undirected
for (int i = c.start[3]; i < c.start[4]; i++) cout << c.to[i] << " ";
\end{lstlisting}
\begin{lstlisting}[style=trace]
adj[3]: 1 2 4
csr.start: 0 2 4 6 9 10 10
csr.to: 1 2 0 3 0 3 1 2 4 3
\end{lstlisting}

\textbf{Pitfalls.} A self-loop $u{-}u$ appears \emph{twice} in \texttt{adj[u]} (it adds 2 to the degree, as the handshake lemma wants); the CSR keeps the input order, so the two representations visit neighbours in the same order and give identical traversals.

\lstinputlisting[style=cpp]{code/lib/graph.hpp}
\cltest{code/tests/test_graph.cpp}{both representations and the readers against a naive build from the edge list, on random multigraphs with self-loops, directed and undirected; CSR neighbour order must equal input order.}{graph 3000}

%============================================================
\clneed
\section{Iterative DFS}
\label{sec:cl-iterdfs}

\textbf{What it does.} A depth-first search that never recurses. It keeps its own stack of vertices and, for each vertex, an index \texttt{it[u]} of the next neighbour to try; ``calling'' \texttt{dfs(v)} is pushing $v$, ``returning'' is popping $u$ when \texttt{it[u]} has run off the end of its list. That one index is what makes the visit order, the parents and the times \emph{identical} to the recursive DFS (the test checks this on 4000 graphs). \texttt{run(adj)} builds a DFS forest over all vertices; \texttt{run(adj, root)} explores from one root.

\textbf{Times.} One clock ticks on entry \emph{and} on exit, so \tin{v} and \tout{v} are distinct numbers in $0..2n-1$ and the intervals $[\tin{v},\tout{v}]$ are nested exactly like the DFS tree: $u$ is an ancestor of $v$ iff $\tin{u}\le\tin{v}$ and $\tout{v}\le\tout{u}$ (\texttt{is\_ancestor}).

\textbf{Complexity.} $\bigO{n+m}$ time, $\bigO{n}$ extra memory; safe on a path of $10^6$ vertices (the test runs one).

\begin{diagrambox}
\begin{tikzpicture}[x=1cm, y=1cm]
    \begin{scope}[clv0/.style={gvsrc}, clv1/.style={gvdone}, clv2/.style={gvdone}, clv3/.style={gvdone},
                  clv4/.style={gvdone}, clv5/.style={gvsrc},
                  cle01/.style={gtree}, cle13/.style={gtree}, cle23/.style={gtree}, cle34/.style={gtree},
                  cle02/.style={gback, -}]
        \pic (L) {clL};
        \foreach \v/\t in {0/{0,9}, 1/{1,8}, 2/{3,4}, 3/{2,7}, 4/{5,6}, 5/{10,11}}
            \node[gdist, below=1pt of L-\v] {[\t]};
    \end{scope}
    \begin{scope}[xshift=7.2cm, yshift=-1cm, x=0.36cm]
        \foreach \t in {0,...,11} \node[font=\tiny\color{darkgray}] at (\t,-0.45) {\t};
        \foreach \v/\a/\b/\y in {0/0/9/2.4, 1/1/8/1.95, 3/2/7/1.5, 2/3/4/1.05, 4/5/6/1.05, 5/10/11/2.4} {
            \draw[primaryblue, line width=3pt] (\a,\y) -- (\b,\y);
            \node[font=\scriptsize\bfseries, anchor=east] at (\a-0.1,\y) {\v};
        }
        \node[font=\scriptsize\color{darkgray}, align=center] at (5.5,3.0) {intervals [tin, tout]: nested = ancestor};
    \end{scope}
\end{tikzpicture}
\figcap{\texttt{IterDFS::run} on $L$ (real output). Thick edges: the DFS forest (roots 0 and 5); the dashed edge $0{-}2$ is a non-tree edge. Vertex 3 is entered from 1, and 2 from 3, so the DFS goes $0\to1\to3\to2$, back to 3, then $3\to4$.}
\end{diagrambox}

\begin{lstlisting}[style=snippet]
IterDFS d;
d.run(g.adj);                    // the whole forest
// d.is_ancestor(1, 4) == true,  d.is_ancestor(2, 4) == false
\end{lstlisting}
\begin{lstlisting}[style=trace]
order: 0 1 3 2 4 5
parent: -1 0 3 1 3 -1
tin: 0 1 3 2 5 10
tout: 9 8 4 7 6 11
is_ancestor(1,4)=1 is_ancestor(2,4)=0
\end{lstlisting}

\textbf{Pitfalls.} The tempting shortcut ``pop $u$, push \emph{all} unvisited neighbours'' is \emph{not} a DFS: it visits vertices in a different order and has no meaningful \tout{}. Use it only when any traversal order is fine (plain reachability).

\lstinputlisting[style=cpp]{code/lib/dfs_iterative.hpp}
\cltest{code/tests/test_dfs_iterative.cpp}{visited set equals Floyd--Warshall reachability; parent arcs exist; parenthesis property of [tin, tout]; no arc leads to a vertex discovered after its tail finished; \texttt{is\_ancestor} equals naive parent climbing; output identical to \texttt{RecDFS}; plus a path of $10^6$ vertices.}{dfs\_iterative 4001}

%============================================================
\clneed
\section{Recursive DFS}
\label{sec:cl-recdfs}

\textbf{What it does.} The textbook DFS, with exactly the outputs of \texttt{IterDFS}. It is shorter and easier to extend (low-links, subtree sizes, DP on trees are written ``after the recursive call''), which is why many chapters show it first.

\textbf{Complexity.} $\bigO{n+m}$ time; recursion depth equals the longest path in the DFS tree, up to $n-1$.

\begin{warningbox}[The depth note]
Every open call is a stack frame. The frames of this \texttt{dfs} are small, yet the program \clfile{code/examples/pitfalls/p1_depth.cpp} (recursive DFS on a path) with the default 2\,MB stack of a MinGW executable on Windows finished for $n=60\,000$ and crashed for $n=70\,000$ with exit code \texttt{0xC00000FD} (stack overflow). Linked with \texttt{-Wl,--stack,268435456} (256\,MB, the Codeforces setting) it finished $n=10^6$. Rule: if the depth can exceed $\approx 5\cdot10^4$ and you do not control the stack size, use \texttt{IterDFS}.
\end{warningbox}

\begin{diagrambox}
\begin{tikzpicture}[x=1cm, y=1cm]
    \foreach \n/\h/\c/\l in {1/1.0/successgreen/{$10^4$}, 2/1.6/successgreen/{$5\cdot10^4$}, 3/1.9/successgreen/{$6\cdot10^4$}, 4/2.2/warningred/{$7\cdot10^4$}, 5/3.1/warningred/{$10^5$}, 6/4.0/warningred/{$10^6$}} {
        \fill[\c!30, draw=\c] (\n*1.3-0.4,0) rectangle (\n*1.3+0.4,\h);
        \node[font=\scriptsize] at (\n*1.3,-0.3) {\l};
    }
    \draw[warningred, dashed, line width=1pt] (0.3,2.05) -- (8.5,2.05) node[right, font=\scriptsize\bfseries, align=left] {2 MB stack\\exhausted};
    \node[font=\scriptsize\color{successgreen}\bfseries] at (2.6,2.4) {finished};
    \node[font=\scriptsize\color{warningred}\bfseries] at (6.5,4.3) {0xC00000FD};
    \node[font=\scriptsize\color{darkgray}] at (4.5,-0.75) {path length $n$ = recursion depth (bar height $\sim\log n$)};
\end{tikzpicture}
\figcap{Measured: recursive DFS on a path of $n$ vertices, default stack, Windows/MinGW.}
\end{diagrambox}

\begin{lstlisting}[style=snippet]
RecDFS r;
r.run(g.adj);                    // same interface and output as IterDFS
\end{lstlisting}
\begin{lstlisting}[style=trace]
rec order: 0 1 3 2 4 5
\end{lstlisting}

\lstinputlisting[style=cpp]{code/lib/dfs_recursive.hpp}
\cltest{code/tests/test_dfs_recursive.cpp}{the same structural checks as for \texttt{IterDFS}, plus ``one DFS tree per connected component''.}{dfs\_recursive 4000}

%============================================================
\clneed
\section{BFS}
\label{sec:cl-bfs}

\textbf{What it does.} Breadth-first search from one source: \texttt{dist[v]} is the fewest number of edges from $s$ (or $-1$), \texttt{par[v]} the vertex that discovered $v$, and \texttt{path(t)} walks the parents back from $t$ and reverses them. The queue is a plain \texttt{vector} with a moving head (no \texttt{std::queue}), which also gives the pop order for free.

\textbf{Complexity.} $\bigO{n+m}$.

\begin{diagrambox}
\begin{tikzpicture}[x=1cm, y=1cm]
    \foreach \d/\x in {0/0, 1/1.3, 2/2.6, 3/3.9} {
        \fill[primaryblue!\the\numexpr 6+5*\d\relax, rounded corners=4pt] (\x-0.5,-1.35) rectangle (\x+0.5,1.35);
        \node[font=\scriptsize\bfseries\color{deepblue}] at (\x,1.6) {dist \d};
    }
    \begin{scope}[clv0/.style={gvsrc}, clv4/.style={gvdst}, clv5/.style={gvdead},
                  cle01/.style={gpath}, cle02/.style={gtree}, cle13/.style={gpath}, cle34/.style={gpath},
                  cle23/.style={gfaded, dashed}]
        \pic (L) {clL};
    \end{scope}
    \node[font=\scriptsize\color{darkgray}] at (5.1,-0.7) {unreachable};
\end{tikzpicture}
\figcap{\texttt{BFS::run(g.adj, 0)} on $L$. The BFS tree is $0{\to}1, 0{\to}2, 1{\to}3, 3{\to}4$; orange: \texttt{path(4)} $=0,1,3,4$. Vertex 3 is discovered from 1 (popped before 2), so the dashed edge $2{-}3$ is unused.}
\end{diagrambox}

\begin{lstlisting}[style=snippet]
BFS b;
b.run(g.adj, 0);
vector<int> p = b.path(4);      // {0, 1, 3, 4};  b.path(5) is empty
\end{lstlisting}
\begin{lstlisting}[style=trace]
dist: 0 1 1 2 3 -1
par: -1 0 0 1 3 -1
path(4): 0 1 3 4
order: 0 1 2 3 4
\end{lstlisting}

\textbf{Pitfalls.} Set \texttt{dist[v]} when $v$ is \emph{pushed}. Appendix~\ref{ch:testing} shows, with a real stress test, how the ``mark when popped'' version produces a wrong distance on a triangle.

\lstinputlisting[style=cpp]{code/lib/bfs.hpp}
\cltest{code/tests/test_bfs.cpp}{\texttt{dist} against Floyd--Warshall on random (multi)graphs with self-loops, directed and undirected; every \texttt{path(t)} has \texttt{dist[t]+1} vertices, starts at $s$, ends at $t$ and uses real edges; the pop order is sorted by distance.}{bfs 5000}

%============================================================
\clneed
\section{Multi-Source BFS}
\label{sec:cl-multibfs}

\textbf{What it does.} The distance from every vertex to the \emph{nearest} of several sources, in one BFS: all sources start in the queue with distance 0. It is the same as one BFS from an imaginary super-source joined to every source. \texttt{src[v]} tells \emph{which} source is nearest (the one whose wave reached $v$ first), \texttt{par[v]} the shortest-path tree.

\textbf{Complexity.} $\bigO{n+m}$, no matter how many sources --- $k$ separate BFS runs would cost $\bigO{k(n+m)}$.

\begin{diagrambox}
\begin{tikzpicture}[x=1cm, y=1cm]
    \begin{scope}[clv0/.style={gvsrc}, clv4/.style={gvsrc}, clv1/.style={gvA}, clv2/.style={gvA}, clv3/.style={gvB},
                  clv5/.style={gvdead}, cle01/.style={gtree}, cle02/.style={gtree}, cle34/.style={gtree, draw=recursionred}]
        \pic (L) {clL};
        \foreach \v/\d in {0/0, 1/1, 2/1, 3/1, 4/0} \node[gdist, above=1pt of L-\v] {\d};
    \end{scope}
    \node[gvsrc, minimum size=6mm, font=\scriptsize\bfseries] (S) at (2.0,2.0) {S};
    \draw[gfaded, dashed, -{Stealth[length=2mm]}] (S) to[bend right=20] (L-0);
    \draw[gfaded, dashed, -{Stealth[length=2mm]}] (S) to[bend left=20] (L-4);
    \node[note, anchor=west] at (2.5,2.0) {imaginary super-source, edges of length 0};
\end{tikzpicture}
\figcap{\texttt{MultiBFS::run(g.adj, \{0, 4\})} on $L$. Blue vertices are nearest to 0, red to 4; vertex 3 is one step from 4 but two from 0.}
\end{diagrambox}

\begin{lstlisting}[style=snippet]
MultiBFS mb;
mb.run(g.adj, {0, 4});
\end{lstlisting}
\begin{lstlisting}[style=trace]
dist: 0 1 1 1 0 -1
src: 0 0 0 4 4 -1
\end{lstlisting}

\lstinputlisting[style=cpp]{code/lib/multi_bfs.hpp}
\cltest{code/tests/test_multi_bfs.cpp}{\texttt{dist[v]} equals the minimum over the sources of the Floyd--Warshall distance (duplicate sources included); \texttt{src[v]} is a source at exactly that distance; \texttt{par} is a shortest-path tree that stays within one source's region.}{multi\_bfs 5000}

%============================================================
\clneed
\section{Grid BFS}
\label{sec:cl-gridbfs}

\textbf{What it does.} BFS directly on a character grid, without building a graph: cell $(r,c)$ is a vertex, walls (\texttt{'\#'} by default) are missing vertices, and the edges are the 4 side moves --- or all 8 king moves with \texttt{eight = true}. Several start cells may be given (multi-source). \texttt{path(r,c)} returns the cells of a shortest route.

\textbf{Complexity.} $\bigO{R\cdot C\cdot k}$ with $k\in\{4,8\}$.

\begin{diagrambox}
\begin{tikzpicture}[x=0.75cm, y=0.75cm]
    \foreach \r/\c in {0/0, 0/1, 0/5, 0/6, 1/0, 2/2, 2/5, 3/4, 3/5} \node[gwall, minimum size=7mm] at (\c,-\r) {};
    \foreach \r/\c/\d in {0/2/2, 0/3/3, 0/4/4, 1/1/0, 1/2/1, 1/3/2, 1/4/3, 1/5/4, 1/6/5, 2/0/2, 2/1/1, 2/3/3, 2/4/4, 2/6/6,
                          3/0/3, 3/1/2, 3/2/3, 3/3/4, 3/6/7}
        \node[gcell, minimum size=7mm] (g\r\c) at (\c,-\r) {\d};
    \node[gcell, minimum size=7mm, fill=basefill, font=\small\bfseries] at (1,-1) {0};
    \node[gcell, minimum size=7mm, fill=answerfill, font=\small\bfseries] at (6,-3) {7};
    \foreach \r/\c/\d in {1/2/1, 1/3/2, 1/4/3, 1/5/4, 1/6/5, 2/6/6} \node[gcell, minimum size=7mm, fill=pathfill] at (\c,-\r) {\d};
    \foreach \a/\b in {11/12, 12/13, 13/14, 14/15, 15/16, 16/26, 26/36} \draw[accentorange, line width=1.2pt, -{Stealth[length=1.6mm]}] (g\a) -- (g\b);
    \node[note, anchor=west, align=left] at (7.6,-1.5) {grid from\\gen.py grid 4 7 30 6;\\start (1,1);\\4 moves: (0,0) unreachable,\\dist(3,6) = 7;\\8 moves: dist(3,6) = 6,\\dist(0,0) = 1};
\end{tikzpicture}
\figcap{\texttt{GridBFS::run(grid, \{\{1,1\}\})} with 4 moves (real output). Dark cells are walls; the shaded cells and arrows are \texttt{path(3,6)}.}
\end{diagrambox}

\begin{lstlisting}[style=snippet]
vector<string> grid = {".#...##", "#......", "..#..#.", "....##."};
GridBFS gb;
gb.run(grid, {{1, 1}});          // 4 moves;  gb.run(grid, {{1, 1}}, true) for 8 moves
auto p = gb.path(3, 6);
\end{lstlisting}
\begin{lstlisting}[style=trace]
path(3,6): (1,1) (1,2) (1,3) (1,4) (1,5) (1,6) (2,6) (3,6)
8-dir dist(3,6)=6 dist(0,0)=1 4-dir dist(0,0)=-1
\end{lstlisting}

\textbf{Pitfalls.} Check the bounds before reading \texttt{g[nr][nc]}. With 8 moves, the step from $(1,1)$ to $(0,0)$ squeezes diagonally between the walls $(0,1)$ and $(1,0)$; if the statement forbids cutting corners, test the two side cells too.

\lstinputlisting[style=cpp]{code/lib/grid_bfs.hpp}
\cltest{code/tests/test_grid_bfs.cpp}{distances against Floyd--Warshall on the explicitly built cell graph (random grids up to $5\times5$, 0--50\% walls, 1--3 starts, 4 and 8 moves); every path step is a legal move onto a free cell.}{grid\_bfs 3000}

%============================================================
\clneed
\section{Connected Components}
\label{sec:cl-components}

\textbf{What it does.} Labels the connected components of an undirected graph: \texttt{comp[v]} in $0..\texttt{count}-1$, numbered in order of their smallest vertex (so the labelling is canonical and easy to compare), and \texttt{size[c]}. One BFS per still-unlabelled vertex.

\textbf{Complexity.} $\bigO{n+m}$ in total: every vertex is labelled once over all the searches.

\begin{diagrambox}
\begin{tikzpicture}
    \begin{scope}[clv0/.style={gvA}, clv1/.style={gvA}, clv2/.style={gvA}, clv3/.style={gvA}, clv4/.style={gvA}, clv5/.style={gvB}]
        \pic (L) {clL};
    \end{scope}
    \node[draw=primaryblue, dashed, rounded corners=6pt, line width=1pt, fit=(L-0)(L-1)(L-2)(L-4), inner sep=5pt,
          label={[font=\scriptsize\bfseries, text=primaryblue]below:component 0, size 5}] {};
    \node[draw=recursionred, dashed, rounded corners=6pt, line width=1pt, fit=(L-5), inner sep=5pt,
          label={[font=\scriptsize\bfseries, text=recursionred]below:component 1}] {};
\end{tikzpicture}
\figcap{\texttt{Components::run(g.adj)} on $L$: two components.}
\end{diagrambox}

\begin{lstlisting}[style=snippet]
Components cc;
cc.run(g.adj);
\end{lstlisting}
\begin{lstlisting}[style=trace]
count=2
comp: 0 0 0 0 0 1
size: 5 1
\end{lstlisting}

\textbf{Pitfalls.} For a \emph{directed} graph this is the wrong tool: add the reverse arcs for weak components, or use SCC (Chapter~19) for strong ones. When edges only arrive (never leave), a DSU (Section~\ref{sec:cl-dsu}) answers ``same component?'' online.

\lstinputlisting[style=cpp]{code/lib/components.hpp}
\cltest{code/tests/test_components.cpp}{same id iff Floyd--Warshall says mutually reachable; ids appear in order of the smallest vertex; sizes match.}{components 5000}

%============================================================
\clneed
\section{Bipartite Check}
\label{sec:cl-bipartite}

\textbf{What it does.} Tries to colour the vertices with two colours so that every edge joins different colours. BFS colours by the parity of the depth. If an edge $u{-}v$ joins two vertices of the same colour, the two tree paths from $u$ and $v$ up to their lowest common ancestor plus the edge itself form an odd cycle --- a \emph{certificate} that no colouring exists. So the answer comes with proof either way: a colouring, or an odd cycle.

\textbf{Complexity.} $\bigO{n+m}$.

\begin{diagrambox}
\begin{tikzpicture}
    \begin{scope}[clv0/.style={gvA}, clv1/.style={gvB}, clv2/.style={gvB}, clv3/.style={gvA}, clv4/.style={gvB}, clv5/.style={gvA}]
        \pic (L) {clL};
        \node[font=\scriptsize\bfseries] at (2.55,-1.55) {$L$: bipartite, colours 0 1 1 0 1 0};
    \end{scope}
    \begin{scope}[xshift=7.6cm, clv0/.style={gvdst}, clv1/.style={gvdst}, clv2/.style={gvdst}, cle01/.style={gpath}, cle02/.style={gpath}]
        \pic (M) {clL};
        \draw[gpath] (M-1) -- (M-2);
        \node[font=\scriptsize\bfseries] at (2.55,-1.55) {$L+\{1{-}2\}$: odd cycle 1, 0, 2};
    \end{scope}
\end{tikzpicture}
\figcap{Both certificates (real output). Left: colour classes blue/red. Right: the extra edge $1{-}2$ joins two vertices at depth 1; their common ancestor 0 closes the triangle.}
\end{diagrambox}

\begin{lstlisting}[style=snippet]
Bipartite bp;
bool ok = bp.run(g.adj);         // true; bp.color is a valid colouring
Graph g2 = g; g2.add_edge(1, 2);
ok = bp.run(g2.adj);             // false; bp.odd_cycle = {1, 0, 2}
\end{lstlisting}
\begin{lstlisting}[style=trace]
L bipartite=1
color: 0 1 1 0 1 0
L+{1-2} bipartite=0
odd_cycle: 1 0 2
L+{3-4 again} bipartite=1
\end{lstlisting}

\textbf{Pitfalls.} Start a BFS from \emph{every} uncoloured vertex (a non-bipartite piece may hide in another component); a self-loop is an odd cycle of length 1 (\texttt{odd\_cycle = \{u\}}); parallel edges change nothing (last line of the output).

\lstinputlisting[style=cpp]{code/lib/bipartite.hpp}
\cltest{code/tests/test_bipartite.cpp}{the yes/no answer against trying all $2^n$ colourings; a ``yes'' must come with a proper colouring, a ``no'' with an odd-length cycle of distinct vertices whose consecutive pairs (and last--first) are edges.}{bipartite 5000}

%============================================================
\clneed
\section{Finding a Cycle}
\label{sec:cl-cycle}

\textbf{What it does.} Returns one cycle, or nothing if there is none.
\begin{itemize}[leftmargin=*]
    \item \textbf{Undirected} (\texttt{find\_cycle\_undirected(n, edges)}): works on the edge list because a cycle in a multigraph is described by its \emph{edges}: two parallel edges form a cycle of length 2, a self-loop one of length 1. The DFS skips only the edge it arrived by (its id), never the parent vertex. Returns the vertices and the ids of the edges between consecutive vertices.
    \item \textbf{Directed} (\texttt{find\_cycle\_directed(adj)}): three colours; an arc into a vertex that is still on the stack (colour 1) closes a cycle, which is read off the parent pointers.
\end{itemize}

\textbf{Complexity.} $\bigO{n+m}$, iterative.

\begin{diagrambox}
\begin{tikzpicture}
    \begin{scope}[cle01/.style={gpath}, cle02/.style={gpath}, cle13/.style={gpath}, cle23/.style={gpath}]
        \pic (L) {clL};
        \node[font=\scriptsize\bfseries, align=center] at (2.55,-1.6) {undirected $L$: verts 2 3 1 0,\\edge ids 3 2 0 1};
    \end{scope}
    \begin{scope}[xshift=7.2cm, yshift=0.2cm]
        \node[gv] (a) at (0,0) {0}; \node[gv] (b) at (1.4,0) {1}; \node[gv] (c) at (2.8,0) {2};
        \draw[ge] (a) -- node[gw] {0} (b);
        \draw[gpath] (b) to[bend left=35] node[gw] {1} (c);
        \draw[gpath] (b) to[bend right=35] node[gw] {2} (c);
        \node[font=\scriptsize\bfseries, align=center] at (1.4,-1.8) {multigraph 0-1, 1-2, 1-2:\\verts 2 1, edge ids 1 2};
    \end{scope}
\end{tikzpicture}
\figcap{Two undirected answers (real output): in $L$ the 4-cycle; in the multigraph the two parallel edges, a cycle of length 2 that a ``skip the parent vertex'' DFS would miss.}
\end{diagrambox}

\begin{diagrambox}
\begin{tikzpicture}[x=1.25cm]
    % gen dag 6 7 4: 5->0 3->0 4->2 3->5 5->4 3->1 3->4, drawn in the Kahn order 3 5 1 0 4 2, plus 2->3
    \foreach \v/\x in {3/0, 5/1, 1/2, 0/3, 4/4, 2/5} \node[gv] (d\v) at (\x,0) {\v};
    \node[gvdst] at (d3) {3}; \node[gvdst] at (d5) {5}; \node[gvdst] at (d4) {4}; \node[gvdst] at (d2) {2};
    \draw[gde] (d3) -- (d5);
    \draw[gde] (d3) to[bend left=35] (d1);
    \draw[gde] (d3) to[bend left=45] (d0);
    \draw[gde] (d5) to[bend right=40] (d0);
    \draw[gde] (d5) to[bend right=45] (d4);
    \draw[gde] (d3) to[bend right=50] (d4);
    \draw[gde] (d4) -- (d2);
    \draw[gpath, -{Stealth[length=2.6mm]}] (d3) -- (d5);
    \draw[gpath, -{Stealth[length=2.6mm]}] (d5) to[bend right=45] (d4);
    \draw[gpath, -{Stealth[length=2.6mm]}] (d4) -- (d2);
    \draw[gpath, -{Stealth[length=2.6mm]}, dashed] (d2) to[bend right=30] node[gw, above] {added} (d3);
\end{tikzpicture}
\figcap{A random DAG (\texttt{gen.py dag 6 7 4}) has no directed cycle: \texttt{find\_cycle\_directed} returns an empty vector. After adding the arc $2\to3$ it returns $2,3,5,4$, i.e.\ $2\to3\to5\to4\to2$ (orange).}
\end{diagrambox}

\begin{lstlisting}[style=snippet]
UndirectedCycle uc = find_cycle_undirected(6, E);              // E = edges of L
vector<int> dc = find_cycle_directed(to_adj(6, D2, true));     // D2 = DAG + {2 -> 3}
\end{lstlisting}
\begin{lstlisting}[style=trace]
verts: 2 3 1 0
ids: 3 2 0 1
multi verts: 2 1
multi ids: 1 2
dag cycle:
dag+2->3 cycle: 2 3 5 4
\end{lstlisting}

\lstinputlisting[style=cpp]{code/lib/cycle.hpp}
\cltest{code/tests/test_cycle.cpp}{undirected: a cycle is found iff $m > n - (\text{number of components})$ (true for multigraphs with loops), the cycle has distinct vertices and edges and edge \texttt{ids[i]} joins \texttt{verts[i]} and \texttt{verts[i+1]}; directed: found iff some vertex reaches itself in the transitive closure, and every arc of the cycle exists.}{cycle 4000}

%============================================================
\clneed
\section{Topological Sort}
\label{sec:cl-topo}

\textbf{What it does.} Orders the vertices of a directed graph so that every arc goes forward, or reports that this is impossible (a cycle) by returning an \emph{empty} vector.
\begin{itemize}[leftmargin=*]
    \item \texttt{topo\_kahn}: count in-degrees, put every vertex of in-degree 0 in a queue, and ``remove'' popped vertices by decrementing their successors' in-degrees. If fewer than $n$ vertices come out, the rest is stuck behind a cycle. Replacing the queue by a min-heap gives the lexicographically smallest order.
    \item \texttt{topo\_dfs}: the reverse of the DFS post-order (a vertex finishes only after everything it reaches has finished); an arc into the current stack is a cycle.
\end{itemize}

\textbf{Complexity.} $\bigO{n+m}$ (Kahn with a heap: $\bigO{m\log n}$).

\begin{diagrambox}
\begin{tikzpicture}[x=1.25cm]
    \foreach \v/\x [count=\i] in {3/0, 5/1, 1/2, 0/3, 4/4, 2/5} {
        \node[gv] (k\v) at (\x,0) {\v};
        \node[font=\tiny\color{darkgray}] at (\x,-0.6) {pos \the\numexpr\i-1\relax};
    }
    \draw[gde] (k3) -- (k5);
    \draw[gde] (k3) to[bend left=35] (k1);
    \draw[gde] (k3) to[bend left=45] (k0);
    \draw[gde] (k5) to[bend right=40] (k0);
    \draw[gde] (k5) to[bend right=45] (k4);
    \draw[gde] (k3) to[bend right=55] (k4);
    \draw[gde] (k4) -- (k2);
    \node[font=\scriptsize\bfseries, anchor=west] at (6.4,0.3) {Kahn: 3 5 1 0 4 2};
    \node[font=\scriptsize\bfseries, anchor=west] at (6.4,-0.3) {DFS: \ \ 3 5 4 2 1 0};
\end{tikzpicture}
\figcap{The DAG of the previous section laid out in the order returned by \texttt{topo\_kahn}: all seven arcs point to the right. \texttt{topo\_dfs} returns a different, equally valid order.}
\end{diagrambox}

\begin{lstlisting}[style=snippet]
vector<int> a = topo_kahn(to_adj(6, D, true));    // D = gen.py dag 6 7 4
vector<int> b = topo_dfs(to_adj(6, D, true));
vector<int> c = topo_kahn(to_adj(6, D2, true));   // with the arc 2 -> 3: empty
\end{lstlisting}
\begin{lstlisting}[style=trace]
kahn: 3 5 1 0 4 2
dfs: 3 5 4 2 1 0
kahn with cycle size=0
\end{lstlisting}

\lstinputlisting[style=cpp]{code/lib/toposort.hpp}
\cltest{code/tests/test_toposort.cpp}{on random DAGs and random digraphs (self-loops and parallel arcs included), both functions return empty iff the transitive closure has a vertex reaching itself, and otherwise a permutation in which every arc points forward.}{toposort 5000}

%============================================================
\clneed
\section{Functional Graphs}
\label{sec:cl-functional}

\textbf{What it does.} In a functional graph every vertex has exactly one outgoing arc $v\to f(v)$. Every component is one directed cycle with trees hanging into it (a ``rho'' shape). \texttt{build(f)} finds all cycles in arc order and, for every vertex, the cycle it falls into (\texttt{cycle\_id}), how many steps it takes to get there (\texttt{dist\_to\_cycle}), the first cycle vertex it hits (\texttt{entry}) and, for cycle vertices, the position in the cycle (\texttt{pos}). With these, $f^{k}(v)$ for large $k$ is one array lookup (formula in the header).

\textbf{Complexity.} $\bigO{n}$; iterative, with three states per vertex: unseen, on the current walk, finished. A walk that bites its own tail has found a new cycle; a walk that hits a finished vertex inherits its data.

\begin{diagrambox}
\begin{tikzpicture}
    \foreach \v [count=\i from 0] in {0,4,7,2,8,1,5,3} \node[gvA] (f\v) at ({90-45*\i}:1.55) {\v};
    \foreach \a/\b in {0/4, 4/7, 7/2, 2/8, 8/1, 1/5, 5/3, 3/0} \draw[gde] (f\a) -- (f\b);
    \node[gv] (f6) at (-2.9,1.55) {6};
    \draw[gde] (f6) -- (f0);
    \node[gdist, above=1pt of f6] {dist 1};
    \node[note, anchor=west, align=left] at (2.6,0) {f = 4 5 8 0 7 3 0 2 1\\(gen.py functional 9 5)\\[2pt]one cycle of length 8:\\0 4 7 2 8 1 5 3\\vertex 6: dist 1, entry 0};
\end{tikzpicture}
\figcap{A random functional graph on 9 vertices (real output): one cycle of length 8 and a single tail vertex.}
\end{diagrambox}

\begin{lstlisting}[style=snippet]
vector<int> f = {4, 5, 8, 0, 7, 3, 0, 2, 1};
FunctionalGraph fg;
fg.build(f);
\end{lstlisting}
\begin{lstlisting}[style=trace]
cycle: 0 4 7 2 8 1 5 3
cycle_id: 0 0 0 0 0 0 0 0 0
dist_to_cycle: 0 0 0 0 0 0 1 0 0
entry: 0 1 2 3 4 5 0 7 8
pos: 0 5 3 7 1 6 -1 2 4
\end{lstlisting}

\lstinputlisting[style=cpp]{code/lib/functional_graph.hpp}
\cltest{code/tests/test_functional_graph.cpp}{against simulation of $f$: on-cycle iff $f^k(v)=v$ for some $k\le n$; distance and entry by walking; cycles in arc order with consistent ids and positions; the $f^k$ formula for all $k\le 2n$.}{functional\_graph 5000}

%============================================================
\clneed
\section{Disjoint Set Union}
\label{sec:cl-dsu}

\textbf{What it does.} Maintains a partition of $0..n-1$ under ``merge the sets of $a$ and $b$'' and ``are $a$ and $b$ together?''. Each set is a rooted tree; \texttt{find} walks to the root and then points every vertex on the way directly at it (\emph{path compression}, done iteratively in two passes); \texttt{unite} hangs the smaller tree below the larger (\emph{union by size}). \texttt{unite} returns whether anything changed --- exactly the test Kruskal needs.

\textbf{Complexity.} $\bigO{\alpha(n)}$ amortised per operation, where $\alpha$ is the inverse Ackermann function ($\le 4$ for any realistic $n$).

\begin{diagrambox}
\begin{tikzpicture}[x=1cm, y=1cm]
    \begin{scope}
        \node[gv] (a0) at (0,0) {0}; \node[gv] (a1) at (0,-1) {1};
        \node[gv] (a2) at (1.3,0) {2}; \node[gv] (a3) at (1.3,-1) {3};
        \node[gv] (a4) at (2.6,0) {4}; \node[gv] (a5) at (3.6,0) {5};
        \draw[gde] (a1) -- (a0); \draw[gde] (a3) -- (a2);
        \node[font=\scriptsize\bfseries] at (1.8,-1.8) {after unite(0,1), unite(2,3)};
    \end{scope}
    \draw[fcarrow] (4.3,-0.5) -- node[lbl, above] {unite(1,3)} (5.6,-0.5);
    \begin{scope}[xshift=6.4cm]
        \node[gvsel] (b0) at (0.65,0.3) {0}; \node[gv] (b1) at (0,-0.8) {1}; \node[gv] (b2) at (1.3,-0.8) {2}; \node[gv] (b3) at (1.3,-1.8) {3};
        \node[gv] (b4) at (2.6,0.3) {4}; \node[gv] (b5) at (3.6,0.3) {5};
        \draw[gde] (b1) -- (b0); \draw[gde] (b2) -- (b0); \draw[gde] (b3) -- (b2);
        \node[font=\scriptsize\bfseries, align=center] at (1.8,-2.6) {sets = 3, size(3) = 4\\unite(0,2) returns false};
    \end{scope}
\end{tikzpicture}
\figcap{DSU on 6 elements (real output): the root of $\{2,3\}$ is hung under the root of $\{0,1\}$ (equal sizes, first argument wins); the next \texttt{find(3)} will point 3 straight at 0.}
\end{diagrambox}

\begin{lstlisting}[style=snippet]
DSU ds(6);
ds.unite(0, 1); ds.unite(2, 3); ds.unite(1, 3);   // true, true, true
ds.unite(0, 2);                                   // false: already together
\end{lstlisting}
\begin{lstlisting}[style=trace]
1110
sets=3 size(3)=4 same(0,3)=1 same(0,4)=0
\end{lstlisting}

\lstinputlisting[style=cpp]{code/lib/dsu.hpp}
\cltest{code/tests/test_dsu.cpp}{random sequences of unite/same on up to 30 elements against an explicit label array that is relabelled on every merge: return values, \texttt{same}, \texttt{size} and \texttt{sets} after every operation.}{dsu 3000}

%============================================================
\clneed
\section{Weighted (Potential) DSU}
\label{sec:cl-wdsu}

\textbf{What it does.} Each element has an unknown number $\mathrm{pot}(v)$ and we learn equations $\mathrm{pot}(b)-\mathrm{pot}(a)=w$ (``$b$ is $w$ taller than $a$'', ``$b$ is $w$ days after $a$''). The DSU stores, for every element, its offset to its parent; \texttt{find} adds the offsets up along the path while compressing it. \texttt{unite(a,b,w)} either merges two sets (choosing the offset between the roots so that the new equation holds) or, inside one set, checks the equation and returns \texttt{false} on a \emph{contradiction}. \texttt{diff(a,b)} answers $\mathrm{pot}(b)-\mathrm{pot}(a)$ for elements of one set.

\textbf{Complexity.} $\bigO{\log n}$ worst case (union by size bounds the depth, so the recursive \texttt{find} is safe), $\bigO{\alpha(n)}$ amortised. Replacing $+$/$-$ by XOR gives the parity DSU.

\begin{diagrambox}
\begin{tikzpicture}[x=1.6cm]
    \node[gv] (p0) at (0,0) {0}; \node[gv] (p1) at (1,0) {1}; \node[gv] (p2) at (2,0) {2}; \node[gv] (p3) at (3.2,0) {3};
    \draw[gde] (p0) -- node[gw, above] {+5} (p1);
    \draw[gde] (p1) -- node[gw, above] {$-2$} (p2);
    \draw[gde, draw=successgreen] (p0) to[bend right=40] node[gw, below] {+3 \yes} (p2);
    \draw[gde, draw=warningred, dashed] (p0) to[bend right=70] node[gw, below] {+4 \no} (p2);
    \node[note, anchor=west, align=left] at (3.6,0) {unite(0,1,5), unite(1,2,$-2$): true\\unite(0,2,3): true (consistent)\\unite(0,2,4): false (contradiction)\\diff(0,2) = 3, diff(2,1) = 2};
\end{tikzpicture}
\figcap{Equations as labelled arcs (real output). The third equation is implied by the first two; the fourth contradicts them and is rejected. Element 3 is unrelated: \texttt{same(0,3)} is false.}
\end{diagrambox}

\begin{lstlisting}[style=snippet]
WeightedDSU w(4);
w.unite(0, 1, 5);      // pot(1) - pot(0) = 5
w.unite(1, 2, -2);     // pot(2) - pot(1) = -2
w.unite(0, 2, 3);      // implied: true
w.unite(0, 2, 4);      // contradiction: false
\end{lstlisting}
\begin{lstlisting}[style=trace]
1110
diff(0,2)=3 diff(2,1)=2 same(0,3)=0
\end{lstlisting}

\lstinputlisting[style=cpp]{code/lib/weighted_dsu.hpp}
\cltest{code/tests/test_weighted_dsu.cpp}{against a naive model with explicit labels and explicit potentials (on a merge, the whole second group is shifted); half of the in-set equations are consistent so that both return values are exercised; \texttt{same} and \texttt{diff} on random pairs after every step.}{weighted\_dsu 3000}

%============================================================
\clneed
\section{0-1 BFS}
\label{sec:cl-01bfs}

\textbf{What it does.} Shortest paths when every weight is 0 or 1, using a double-ended queue: relaxing along a 0-edge pushes to the \emph{front} (same distance), along a 1-edge to the \emph{back} (one more). The deque then always holds distances $d$ followed by $d+1$ --- the order Dijkstra would use, without a heap.

\textbf{Complexity.} $\bigO{n+m}$.

\begin{diagrambox}
\begin{tikzpicture}[x=1.3cm]
    \node[gvsrc] (z0) at (0,0) {0}; \node[gv] (z1) at (1.3,0.9) {1}; \node[gv] (z2) at (2.6,0) {2};
    \node[gv] (z3) at (1.3,-0.9) {3}; \node[gv] (z4) at (3.9,0.9) {4};
    \draw[ge] (z0) -- node[gw] {1} (z1); \draw[gpath] (z1) -- node[gw] {0} (z2);
    \draw[ge] (z0) -- node[gw] {1} (z3); \draw[ge] (z3) -- node[gw] {1} (z2);
    \draw[ge] (z2) -- node[gw] {1} (z4); \draw[ge] (z1) -- node[gw] {1} (z4);
    \foreach \v/\d in {0/0, 1/1, 2/1, 3/1, 4/2} \node[gdist, above right=-1pt of z\v] {\d};
    \begin{scope}[xshift=6.6cm, yshift=0.2cm]
        \foreach \x/\v/\s in {0/2/basecell, 1/3/cell} \node[\s, minimum width=7mm] at (\x*0.75,0) {\v};
        \node[font=\scriptsize, anchor=east] at (-0.45,0) {front};
        \node[font=\scriptsize, anchor=west] at (1.2,0) {back};
        \node[note, align=center] at (0.4,-0.9) {after popping 1: 2 (dist 1, via a 0-edge)\\went to the FRONT, 3 waits at the back};
    \end{scope}
\end{tikzpicture}
\figcap{\texttt{ZeroOneBFS} from 0 (real output: \texttt{dist = 0 1 1 1 2}). Vertex 2 is reached through the 0-edge $1{-}2$ at distance 1, not through 3 at distance 2.}
\end{diagrambox}

\begin{lstlisting}[style=snippet]
vector<vector<pair<int,int>>> z(5);    // (neighbour, weight 0 or 1), both directions
// edges: 0-1 (1), 1-2 (0), 0-3 (1), 3-2 (1), 2-4 (1), 1-4 (1)
ZeroOneBFS zb;
zb.run(z, 0);
\end{lstlisting}
\begin{lstlisting}[style=trace]
dist: 0 1 1 1 2
\end{lstlisting}

\textbf{Pitfalls.} A vertex may be in the deque twice (first via a 1-edge, then improved via a 0-edge); the \texttt{done} flag skips the stale copy.

\lstinputlisting[style=cpp]{code/lib/zero_one_bfs.hpp}
\cltest{code/tests/test_zero_one_bfs.cpp}{distances against exhaustive search over all simple paths (multigraphs with loops, directed and undirected, weights 0/1); every \texttt{par[t]} lies on a shortest path.}{zero\_one\_bfs 4000}

%============================================================
\clneed
\section{Dijkstra}
\label{sec:cl-dijkstra}

\textbf{What it does.} Single-source shortest paths for \emph{non-negative} weights. The \emph{lazy} version never decreases a key: whenever a distance improves it pushes a new pair (distance, vertex), and a popped pair whose distance is no longer current is skipped. \texttt{path(t)} reconstructs a shortest path.

\textbf{Complexity.} $\bigO{(n+m)\log m}$; the heap holds at most $m+1$ pairs.

\begin{diagrambox}
\begin{tikzpicture}
    \begin{scope}[clw0/.style={gvsrc}, clw2/.style={gvdst}, clf30/.style={gpath}, clf43/.style={gpath}, clf42/.style={gpath},
                  clf01/.style={gtree}]
        \pic (W) {clW};
        \foreach \v/\d in {0/0, 1/5, 2/12, 3/4, 4/6} \node[gdist, above=1pt of W-\v] {\d};
    \end{scope}
    \node[note, anchor=west, align=left] at (5.2,0) {W: gen.py connected 5 6 7\\--weights 1 9\\[2pt]path(2) = 0 3 4 2,\\length 4 + 2 + 6 = 12};
\end{tikzpicture}
\figcap{\texttt{Dijkstra::run(wadj, 0)} on the weighted sample $W$ (real output). Thick edges form the shortest-path tree; orange is \texttt{path(2)}. The direct edge $0{-}4$ (9) loses to $0{-}3{-}4$ (6).}
\end{diagrambox}

\begin{lstlisting}[style=snippet]
vector<Edge> W = {{0,1,5}, {0,4,9}, {4,1,3}, {4,3,2}, {3,0,4}, {4,2,6}};
auto wadj = to_wadj(5, W, false);       // from gen.hpp, or build a WGraph
Dijkstra dj;
dj.run(wadj, 0);
\end{lstlisting}
\begin{lstlisting}[style=trace]
dijkstra dist: 0 5 12 4 6
par: -1 0 4 0 3
path(2): 0 3 4 2
\end{lstlisting}

\lstinputlisting[style=cpp]{code/lib/dijkstra.hpp}
\cltest{code/tests/test_dijkstra.cpp}{against a naive Bellman--Ford \emph{and} exhaustive simple paths, with weights in $[0,3]$ and in $[0,10^9]$ (so sums exceed 32 bits); every \texttt{path(t)} uses real edges and has exactly the reported length.}{dijkstra 5000}

%============================================================
\clneed
\section{Bellman--Ford}
\label{sec:cl-bf}

\textbf{What it does.} Shortest paths when weights may be negative. $n-1$ rounds of relaxing every arc give the correct distances if no negative cycle is reachable. The library then runs $n$ more rounds in which any vertex that can still improve, or is reachable from such a vertex, gets $-\infty$: its distance is unbounded below. \texttt{run} returns whether a negative cycle is reachable from $s$. \texttt{find\_negative\_cycle} looks for a negative cycle \emph{anywhere} (all vertices start at distance 0, as if a virtual source reached them) and returns its arcs in order.

\textbf{Complexity.} $\bigO{nm}$.

\begin{diagrambox}
\begin{tikzpicture}[x=1.4cm]
    \node[gvsrc] (b0) at (0,0) {0}; \node[gv] (b1) at (1.3,0.8) {1}; \node[gv] (b2) at (1.3,-0.8) {2};
    \node[gvdead] (b3) at (2.7,0.8) {3}; \node[gvdead] (b4) at (4.1,0.8) {4}; \node[gvdead] (b5) at (5.5,0.8) {5};
    \node[gv] (b6) at (-1.4,0) {6};
    \draw[gde] (b0) -- node[gw] {4} (b1); \draw[gde] (b0) -- node[gw] {2} (b2); \draw[gde] (b2) -- node[gw] {$-1$} (b1);
    \draw[gde] (b1) -- node[gw] {3} (b3);
    \draw[gpath, -{Stealth[length=2.6mm]}] (b3) to[bend left=30] node[gw] {1} (b4);
    \draw[gpath, -{Stealth[length=2.6mm]}] (b4) to[bend left=30] node[gw] {$-2$} (b3);
    \draw[gde] (b4) -- node[gw] {1} (b5); \draw[gde] (b6) -- node[gw] {1} (b0);
    \foreach \v/\d in {0/0, 1/1, 2/2} \node[gdist, below=1pt of b\v] {\d};
    \foreach \v in {3,4,5} \node[gdist, below=1pt of b\v, text=warningred] {$-\infty$};
    \node[gdist, below=1pt of b6, text=darkgray] {INF};
\end{tikzpicture}
\figcap{Bellman--Ford from 0 (real output). The cycle $3\to4\to3$ has weight $1-2=-1$; 3 and 4 lie on it and 5 is behind it, so all three are $-\infty$. Vertex 6 reaches 0 but is not reachable from it. \texttt{find\_negative\_cycle} returns the arcs 5, 4 ($4\to3$, $3\to4$).}
\end{diagrambox}

\begin{lstlisting}[style=snippet]
vector<Edge> B = {{0,1,4}, {0,2,2}, {2,1,-1}, {1,3,3}, {3,4,1}, {4,3,-2}, {4,5,1}, {6,0,1}};
BellmanFord bf;
bool neg = bf.run(7, B, 0);              // true
vector<int> cyc = find_negative_cycle(7, B);
\end{lstlisting}
\begin{lstlisting}[style=trace]
neg=1
dist: 0 1 2 -INF -INF -INF INF
neg cycle edges: 5 4
\end{lstlisting}

\textbf{Pitfalls.} The vertex relaxed in round $n$ proves a negative cycle exists but may only be \emph{behind} it; walk $n$ parent steps first to land on the cycle. For an undirected graph, a single negative edge is already a negative cycle ($u\to v\to u$).

\lstinputlisting[style=cpp]{code/lib/bellman_ford.hpp}
\cltest{code/tests/test_bellman_ford.cpp}{against exhaustive enumeration of simple paths and simple cycles on random digraphs with negative weights, loops and parallel arcs: $-\infty$ exactly for vertices reachable from $s$ through a vertex on a negative simple cycle, otherwise the shortest simple path; \texttt{find\_negative\_cycle} non-empty iff any negative cycle exists, arcs chain up and sum to a negative number.}{bellman\_ford 5000}

%============================================================
\clneed
\section{Floyd--Warshall}
\label{sec:cl-floyd}

\textbf{What it does.} All-pairs shortest paths on a distance matrix. After the outer iteration $k$, \texttt{d[i][j]} is the shortest path from $i$ to $j$ that uses only $0..k$ as intermediate vertices; after the last one, all paths. Negative weights are fine; a negative \texttt{d[i][i]} afterwards means $i$ lies on (or returns through) a negative cycle.

\textbf{Complexity.} $\bigO{n^3}$ time, $\bigO{n^2}$ memory: $n\le 500$ is comfortable.

\begin{diagrambox}
\begin{tikzpicture}
    \begin{scope}[scale=0.85, transform shape]
        \pic (W) {clW};
    \end{scope}
    \begin{scope}[xshift=5.6cm, yshift=1.6cm, x=0.8cm, y=0.65cm]
        \foreach \j in {0,...,4} \node[hdr, minimum width=7mm, minimum height=5mm] at (\j+1,0) {\j};
        \foreach \i/\row in {0/{0,5,12,4,6}, 1/{5,0,9,5,3}, 2/{12,9,0,8,6}, 3/{4,5,8,0,2}, 4/{6,3,6,2,0}} {
            \node[hdr, minimum width=7mm, minimum height=5mm] at (0,-\i-1) {\i};
            \foreach \x [count=\j from 1] in \row {
                \ifnum\x=0 \node[basecell, minimum width=7mm, minimum height=5mm] at (\j,-\i-1) {\x};
                \else \node[cell, minimum width=7mm, minimum height=5mm] at (\j,-\i-1) {\x}; \fi
            }
        }
        \node[anscell, minimum width=7mm, minimum height=5mm] at (3,-1) {12};
    \end{scope}
\end{tikzpicture}
\figcap{\texttt{floyd\_warshall} on $W$ (real output). Row 0 equals Dijkstra from 0; \texttt{has\_negative\_cycle} is false.}
\end{diagrambox}

\begin{lstlisting}[style=snippet]
auto M = make_matrix(5);                 // 0 on the diagonal, FW_INF elsewhere
for (auto& e : W) {                      // keep the MIN of parallel edges
    M[e.u][e.v] = min(M[e.u][e.v], e.w);
    M[e.v][e.u] = min(M[e.v][e.u], e.w);
}
floyd_warshall(M);
\end{lstlisting}

\textbf{Pitfalls.} The loop over $k$ must be the \emph{outermost}; never add to \texttt{FW\_INF} (the guard does it); with negative cycles the values can fall exponentially, so they are clamped at $-$\texttt{FW\_INF}.

\lstinputlisting[style=cpp]{code/lib/floyd_warshall.hpp}
\cltest{code/tests/test_floyd_warshall.cpp}{negative-cycle detection against simple-cycle enumeration; when there is no negative cycle, every entry against exhaustive simple paths (directed graphs with negative weights, undirected with non-negative ones).}{floyd\_warshall 5000}

%============================================================
\clneed
\section{Kruskal}
\label{sec:cl-kruskal}

\textbf{What it does.} A minimum spanning \emph{forest}: sort the edges by weight and take each edge that joins two different trees, tested with a (local, minimal) DSU. Returns the total weight, the indices of the chosen edges and the number of trees; \texttt{components == 1} means it is a spanning tree.

\textbf{Complexity.} $\bigO{m\log m}$ for the sort; the DSU part is almost linear.

\begin{diagrambox}
\begin{tikzpicture}
    \begin{scope}[clf43/.style={gtree}, clf41/.style={gtree}, clf30/.style={gtree}, clf42/.style={gtree},
                  clf01/.style={gfaded}, clf04/.style={gfaded}]
        \pic (W) {clW};
    \end{scope}
    \node[note, anchor=west, align=left] at (5.2,0.1) {sorted: 2 (4-3), 3 (4-1), 4 (3-0),\\5 (0-1) \no{} cycle, 6 (4-2), 9 (0-4) \no\\[2pt]chosen edge ids 3 2 4 5\\weight 2 + 3 + 4 + 6 = 15};
\end{tikzpicture}
\figcap{Kruskal and Prim on $W$ (real output): the same tree of weight 15. The edges of weight 5 and 9 would close cycles.}
\end{diagrambox}

\begin{lstlisting}[style=snippet]
MSTResult k = kruskal(5, W);
\end{lstlisting}
\begin{lstlisting}[style=trace]
weight=15 comps=1
chosen: 3 2 4 5
\end{lstlisting}

\lstinputlisting[style=cpp]{code/lib/kruskal.hpp}
\cltest{code/tests/test_kruskal.cpp}{the weight against the minimum over all edge subsets that form a spanning forest (up to 11 edges, negative weights, loops and parallel edges, disconnected graphs); the chosen edges are acyclic, there are $n-\text{components}$ of them, and they sum to the reported weight.}{kruskal 4000}

%============================================================
\clneed
\section{Prim}
\label{sec:cl-prim}

\textbf{What it does.} Grows one tree from a start vertex, always adding the cheapest edge that leaves it (lazy heap: stale entries for vertices already in the tree are skipped). It restarts from every vertex not yet reached, so on a disconnected graph it also returns a minimum spanning forest. \texttt{parent[v]} is the tree neighbour through which $v$ joined.

\textbf{Complexity.} $\bigO{m\log m}$. On a dense graph ($m\approx n^2/2$) the array version without a heap, $\bigO{n^2}$, is faster.

\begin{lstlisting}[style=snippet]
PrimResult pr = prim(wadj);
\end{lstlisting}
\begin{lstlisting}[style=trace]
prim weight=15 comps=1
parent: -1 4 4 0 3
\end{lstlisting}

The parents describe the same tree as Kruskal's chosen edges: $1{-}4$, $2{-}4$, $3{-}0$, $4{-}3$.

\textbf{Pitfalls.} In Prim a vertex joins the tree when it is \emph{popped} (its cheapest connection is final only then) --- the opposite of BFS, where a vertex is final when pushed.

\lstinputlisting[style=cpp]{code/lib/prim.hpp}
\cltest{code/tests/test_prim.cpp}{the weight and the number of trees against the exhaustive minimum spanning forest; every parent edge exists, the parent pointers contain no cycle, one root per tree.}{prim 4000}

%============================================================
\clneed
\section{LCA by Binary Lifting}
\label{sec:cl-lca}

\textbf{What it does.} Preprocesses a rooted tree so that the lowest common ancestor of any two vertices is found in $\bigO{\log n}$. \texttt{up[j][v]} is the $2^j$-th ancestor of $v$ (the root is its own ancestor); the $k$-th ancestor is the composition of the jumps of the binary digits of $k$. For \texttt{lca(u,v)}: lift the deeper vertex to the depth of the other, then lift both together by decreasing powers of two as long as their ancestors differ; the parent of where they stop is the answer. The tree is explored with a BFS, so a path-shaped tree of $10^6$ vertices is no problem.

\textbf{Complexity.} $\bigO{n\log n}$ build and memory; $\bigO{\log n}$ per query.

\begin{diagrambox}
\begin{tikzpicture}[x=1.2cm, y=0.95cm]
    % gen.py tree 8 1, rooted at 0
    \node[gvsrc] (t0) at (0,0) {0};
    \node[gv] (t5) at (0,-1) {5};
    \node[gvdst] (t3) at (0,-2) {3};
    \node[gv] (t4) at (-1.3,-3) {4}; \node[gv] (t6) at (0,-3) {6}; \node[gv] (t7) at (1.3,-3) {7};
    \node[gvcur] (t2) at (-1.3,-4) {2}; \node[gvcur] (t1) at (0,-4) {1};
    \draw[gtree] (t0) -- (t5); \draw[gtree] (t5) -- (t3); \draw[ge] (t3) -- (t7);
    \draw[gpath] (t3) -- (t4); \draw[gpath] (t4) -- (t2); \draw[gpath] (t3) -- (t6); \draw[gpath] (t6) -- (t1);
    \foreach \y/\d in {0/0, -1/1, -2/2, -3/3, -4/4} \node[font=\scriptsize\color{darkgray}] at (-2.4,\y) {depth \d};
    \draw[back, dashed] (t2.east) to[bend right=50] node[lbl, right] {up[1][2] = 3} (t3.west);
    \node[note, anchor=west, align=left] at (2.2,-2) {lca(2, 1) = 3, dist(2, 1) = 4\\lca(2, 6) = 3,\ lca(2, 4) = 4\\kth\_ancestor(2, 3) = 5\\kth\_ancestor(2, 5) = $-1$ (too high)\\LOG = 4 \ ($2^4 > 8$)};
\end{tikzpicture}
\figcap{A random tree (\texttt{gen.py tree 8 1}) rooted at 0 (real output). Orange: the path between 2 and 1 through their LCA 3. One jump of length $2^1$ takes 2 straight to 3.}
\end{diagrambox}

\begin{lstlisting}[style=snippet]
vector<pair<int,int>> T = {{3,4}, {1,6}, {5,3}, {3,6}, {0,5}, {2,4}, {7,3}};
LCA L;
L.build(to_adj(8, T, false), 0);
\end{lstlisting}
\begin{lstlisting}[style=trace]
depth: 0 4 4 2 3 1 3 3
up0: 0 6 4 5 3 0 3 3
lca(2,6)=3 lca(1,7)=3 lca(2,4)=4 dist(2,1)=4 kth(2,3)=5 kth(2,5)=-1 LOG=4
\end{lstlisting}

\lstinputlisting[style=cpp]{code/lib/lca.hpp}
\cltest{code/tests/test_lca.cpp}{\texttt{lca}, \texttt{dist} and \texttt{kth\_ancestor} against naive climbing one parent at a time, on random trees (up to 200 vertices) with a random root, 60 queries each.}{lca 3000}

%============================================================
\clneed
\section{Generators}
\label{sec:cl-gen}

The tests are only as good as their inputs. \clfile{code/gen/gen.hpp} (C++) and \clfile{code/gen/gen.py} (Python) contain the same seeded generators: random labelled trees, connected graphs, general graphs with or without self-loops and multi-edges, directed versions, DAGs, functional graphs, grids with walls, and random weights for any edge list. Vertex labels are shuffled and undirected edges get a random orientation, so code that silently relies on ``vertex 0 is the root'' or ``edges come sorted'' is caught. Appendix~\ref{ch:testing} shows what the generators produce and how to use them with \clfile{tools/stress.py}.

\lstinputlisting[style=cpp]{code/gen/gen.hpp}

%============================================================
\clneed
\section{Running the Tests}
\label{sec:cl-run}

\texttt{python tools/test\_library.py} compiles every \clfile{code/tests/test_*.cpp} with \texttt{g++ -std=c++14 -O2}, runs it and prints one line per component; \texttt{python tools/test\_library.py dsu bfs} runs only the matching tests. The exit code is 1 if anything fails, so the script can guard a commit. A real run on the author's machine (g++ 6.3):

\begin{lstlisting}[style=trace]
OK bellman_ford 5000                                       8.2s
OK bfs 5000                                                5.5s
OK bipartite 5000                                          5.2s
OK components 5000                                         4.9s
OK cycle 4000                                              6.1s
OK dfs_iterative 4001                                      5.5s
OK dfs_recursive 4000                                      4.8s
OK dijkstra 5000                                           6.0s
OK dsu 3000                                                6.3s
OK floyd_warshall 5000                                     5.5s
OK functional_graph 5000                                   5.0s
OK graph 3000                                              6.0s
OK grid_bfs 3000                                           4.7s
OK kruskal 4000                                            5.5s
OK lca 3000                                                4.9s
OK multi_bfs 5000                                          4.7s
OK prim 4000                                               1.4s
OK toposort 5000                                           1.2s
OK weighted_dsu 3000                                       1.6s
OK zero_one_bfs 4000                                       3.7s

20/20 library tests passed.
\end{lstlisting}

(The times include compilation.) When a test fails it prints the component, the exact input, the property that broke and the line of the check, and stops. Here is the real output of \clfile{test_bfs.cpp} run against a deliberately broken copy of \clfile{bfs.hpp} that marks vertices when they are popped:

\begin{lstlisting}[style=trace]
FAIL bfs
  input: n=6 m=10 edges: 0-1 0-3 4-3 2-4 3-3 0-0 5-5 3-3 2-1 2-4 undirected s=1
  dist[4] = 3, expected 2  [b.dist[t] == D[s][t]]  (line 21)
\end{lstlisting}

Because every generator is seeded, the same failure reappears on every run until the bug is fixed.

\begin{checklistbox}[Before adding a component to the library]
\begin{itemize}[leftmargin=*]
    \item One header, self-contained, 0-indexed, comment block with what / complexity / indexing / pitfalls.
    \item No recursion deeper than $\bigO{\log n}$, or a depth note as in \clfile{dfs_recursive.hpp}.
    \item Every array is re-\texttt{assign}ed at the start of \texttt{run}/\texttt{build} (multi-test safe).
    \item A test in \clfile{code/tests/} against a brute force that is \emph{obviously} correct, on thousands of random inputs including self-loops, multi-edges, disconnected graphs and $n=1$.
    \item The test returns a \emph{certificate} check whenever the algorithm can return one (a path, a cycle, a colouring, an order) --- not only the yes/no answer.
    \item \texttt{python tools/test\_library.py} prints \texttt{20/20} (or $N/N$) before the change is used in a chapter.
\end{itemize}
\end{checklistbox}
